\documentclass[11pt]{article}
\usepackage{amsmath,amssymb}
\title{A fresh calibration does not identify an outstanding navigation target}
\date{26 November 2026}
\begin{document}
\maketitle
\paragraph{Calibration, not a map-version claim.}
Let $T$ map coordinates in one stable robot-local gauge to the specified
world gauge. A calibration update revises the estimate of $T$; it does not
rename the physical world, change the map origin, or establish that the
new estimate is correct. The pinned adapter keys this transform by
\texttt{map\_id}. This study does not assert that a VDA 5050 map activation
uses this API or that versioned maps may have different origins.

\paragraph{What the source replay separates.}
The unchanged dispatch method looks up the transform and constructs a
\texttt{NodeRequest} by value before handing it to navigation. The
unchanged setter replaces the transform under its map key. It does not
change the handed request. Subsequent dispatch calls decline to send a
replacement while the order is marked executing. A later native
position update uses the current transform, not the one used for the
outstanding request. The serial fixture executes those method bodies
with an inert callback that captures but does not finish the request.

Two runs start with world estimate $x=u+1$ for local coordinate $u$.
Both then use $x=u$. The same released target $x=4$ is dispatched on
opposite sides of the update. They agree on the accepted order, current
transform and position snapshot, including after resampling. The
already-issued local requests are $u=3$ and $u=4$. If the corrected
identity is accurate in the actuation gauge and the consumer retains the
numeric request, the two now target different world points. Delivery of a
request is the observed cut; motion and arrival are not observed.
This example uses sequential calls only, no purported data race.
The source's existing navigation test shows that the transform is used
when dispatch constructs the request; it does not cover this handoff.

\paragraph{An endpoint revalidation law.}
For an outstanding point request $p$, record the local target $u$ that
was actually issued or enough information to derive it, such as the
dispatch transform $T_d$ with $u=T_d^{-1}p$. If $T_c$ is verified and
applies to the same physical local gauge, the existing request meets a
closed Euclidean world-position tolerance $\epsilon$ exactly when
\[
  \left\|T_c u-p\right\|_2\leq\epsilon.
\]
The target in the all-current calculation $T_c(T_c^{-1}p)$ is always
$p$ and is therefore not an evaluation of the old request. If explicit
heading is present, also test its circular distance against its
separate heading tolerance. If no heading was requested, there is no
such obligation. A changed calibration is not itself sufficient to
invalidate an old request: rotation about the target leaves a
position-only request at the same point but can change the required yaw.
The replay checks both outcomes and all 162 integer translation cases
over its declared grid.

If the same transform pair applies to outstanding position-only targets
in a filled, bounded polytope $P$, testing its vertices exactly checks
the whole envelope. To see this, write $T_cT_d^{-1}(p)-p=Ap+b$.
For $p=\sum_i\lambda_i v_i$ in $P$, triangle inequality bounds its
norm by $\sum_i\lambda_i\|Av_i+b\|$, hence by the largest vertex
error, which is attained at a vertex. With a common yaw tolerance the
relative rotation contributes a single circular error for all explicit
yaws. This envelope check can overestimate the error for a smaller discrete set if
the envelope includes extra positions outside its convex hull.

\paragraph{Recovery interface and limits.}
The useful witness is a target actually handed to an actuator, tagged to
the execution instance and its local coordinate gauge. A planner that
separately retains that target and a verified current transform gets the
same endpoint check. Retiring every old request after recalibration
is conservative even for some requests that would be unchanged. Neither
a fresh pose nor the geometric check establishes that a route is clear,
that another component has physically vacated a shared cell, or that
a controller can revoke an executing request. If the current transform
is uncertain, the point and heading tests must be valid for its entire
admissible set; if the local gauge has changed, an actual bridge between
gauges is required. The fixture neither represents that change nor
executes an MQTT stack, controller motion, full map activation, or a
mission acceptance test. No inventory entry is discharged.
\end{document}
